\documentclass{beamer}
\usepackage{amsmath}
\usepackage{amssymb}

% Choose a theme
\usetheme{Madrid}
\usecolortheme{default}
\usefonttheme{serif}

% --- CUSTOM FOOTER ---
% Redefine the footline to remove author/institute and keep title/page number
\makeatletter
\setbeamertemplate{footline}
{
  \leavevmode%
  \hbox{%
  \begin{beamercolorbox}[wd=.5\paperwidth,ht=2.25ex,dp=1ex,center]{title in head/foot}%
    \usebeamerfont{title in head/foot}\insertshorttitle
  \end{beamercolorbox}%
  \begin{beamercolorbox}[wd=.5\paperwidth,ht=2.25ex,dp=1ex,right]{date in head/foot}%
    \usebeamerfont{date in head/foot}\insertframenumber{} / \inserttotalframenumber\hspace*{2ex}
  \end{beamercolorbox}}%
  \vskip0pt%
}
\makeatother
% --- END CUSTOM FOOTER ---

\title[Lecture 21: Equilibrium figures]{Lecture 21: Equilibrium figures}
\author{David Al-Attar}
\institute{Michaelmas term 2024}
\date{}

\begin{document}

% --- INTRODUCTION ---

\begin{frame}
    \titlepage
\end{frame}

\begin{frame}{Outline and Motivation}
    \begin{itemize}
        \item In the last lecture, we showed that the equilibrium state of a rotating, self-gravitating body like the Earth cannot be stress-free.
        \pause
        \item Today, we will investigate these equilibrium equations in more detail.
        \pause
        \item We will see how the requirement of \textbf{hydrostatic equilibrium} places strong constraints on the shape and internal structure of a planet.
        \pause
        \item This will lead to Clairaut's theory, which explains why slowly rotating planets are \textbf{oblate spheroids}.
    \end{itemize}
\end{frame}

% --- HYDROSTATIC EQUILIBRIUM ---

\begin{frame}{The Equilibrium Equations}
    The general equilibrium equation for a steadily rotating, self-gravitating body is a balance of forces:
    \begin{alertblock}{Equilibrium Force Balance}
        $$ \underbrace{\rho\epsilon_{ijk}\epsilon_{klm}\Omega_{j}\Omega_{l}x_{m}}_{\text{Centrifugal Force}} - \underbrace{\frac{\partial\sigma_{ij}}{\partial x_{j}}}_{\text{Stress Divergence}} - \underbrace{\rho g_{i}}_{\text{Gravitational Force}} = 0 $$
    \end{alertblock}
    \pause
    \begin{itemize}
        \item This equation is general and holds for any material. Over geological timescales, planets behave like fluids, so we cannot assume they are purely elastic.
        \pause
        \item For a general stress tensor $\sigma_{ij}$, this is an \textbf{underdetermined} problem (3 equations for 6 unknown stress components).
    \end{itemize}
\end{frame}

\begin{frame}{Hydrostatic Equilibrium (1/3): The Approximation}
    \begin{itemize}
        \item In a fluid that is stationary or flowing very slowly (like the Earth's mantle over long timescales), the stress tensor is isotropic and is described by pressure alone.
    \end{itemize}
    \pause
    \begin{block}{The Hydrostatic Stress Tensor}
        $$ \sigma_{ij} = -p\,\delta_{ij} $$
    \end{block}
    \pause
    \begin{itemize}
        \item Assuming the planet is in \textbf{hydrostatic equilibrium} means we assume this form for the stress.
        \pause
        \item The stresses associated with mantle convection are much smaller than the hydrostatic pressure, so this is an excellent approximation for the Earth's global shape.
    \end{itemize}
\end{frame}

\begin{frame}{Hydrostatic Equilibrium (2/3): The Governing Equation}
    Substituting $\sigma_{ij} = -p\,\delta_{ij}$ into the equilibrium equation gives:
    $$ \rho\epsilon_{ijk}\epsilon_{klm}\Omega_{j}\Omega_{l}x_{m} + \frac{\partial p}{\partial x_{i}} - \rho g_{i} = 0 $$
    \pause
    We can simplify this by writing the forces as gradients of potentials.
    \begin{itemize}
        \item Gravitational force: $\mathbf{g} = -\nabla\phi$
        \item Centrifugal force: $\rho \boldsymbol{\Omega} \times (\boldsymbol{\Omega} \times \mathbf{x}) = \rho \nabla\psi$
    \end{itemize}
    \pause
    \begin{alertblock}{Hydrostatic Equilibrium Equation}
        We introduce the \textbf{gravity potential} $\gamma = \phi + \psi$. The equation then takes the simple form:
        $$ \nabla p = -\rho \nabla\gamma $$
    \end{alertblock}
\end{frame}

\begin{frame}{Hydrostatic Equilibrium (3/3): Consequences}
    The equation $\nabla p = -\rho \nabla\gamma$ has profound consequences.
    \begin{itemize}
        \item It implies that the vector $\nabla p$ is parallel to the vector $\nabla\gamma$. Therefore, surfaces of constant pressure (\textbf{isobars}) must coincide with surfaces of constant gravity potential (\textbf{equipotentials}).
        \pause
        \item Since the planet's surface is a surface of zero pressure, the \textbf{surface of a planet in hydrostatic equilibrium must be an equipotential surface}.
        \pause
        \item By taking the curl of the equation ($\nabla \times (\nabla p / \rho) = \mathbf{0}$), we find:
        $$ \nabla\rho \times \nabla\gamma = \mathbf{0} $$
        This implies that surfaces of constant density must also coincide with the equipotential surfaces.
    \end{itemize}
    \begin{alertblock}{Conclusion}
        For a planet to be in hydrostatic equilibrium, its surfaces of constant pressure, density, and gravity potential must all coincide.
    \end{alertblock}
\end{frame}

% --- CLAIRAUT'S THEORY ---

\begin{frame}{Slowly Rotating Planets (1/3): Perturbation Setup}
    \begin{itemize}
        \item For a non-rotating planet, the only hydrostatic solution is a spherically symmetric one.
        \pause
        \item For the Earth, the centrifugal potential is very small compared to the gravitational potential ($\sim 1/300$). This suggests we can treat rotation as a small perturbation.
        \pause
        \item We expand all quantities as a spherical zeroth-order term plus a small first-order perturbation (denoted by subscript 1).
    \end{itemize}
    \begin{block}{Perturbation Expansions}
        \begin{align*}
            \rho &= \rho_0(r) + s\rho_1(\mathbf{x}) \\
            p &= p_0(r) + sp_1(\mathbf{x}) \\
            \phi &= \phi_0(r) + s\phi_1(\mathbf{x}) \\
            \psi &= s\psi_1(\mathbf{x})
        \end{align*}
        The shape is also perturbed: $r = b + s h_1(\theta, \varphi)$.
    \end{block}
\end{frame}

\begin{frame}{Slowly Rotating Planets (2/3): The Linearised Equation}
    Substituting the expansions into the hydrostatic equilibrium equation and keeping only the first-order terms gives the linearised equation:
    \begin{alertblock}{First-Order Hydrostatic Equilibrium}
        $$ \nabla p_{1} + \rho_{0}\nabla\gamma_{1} + \rho_{1}g_{0}\hat{\mathbf{r}} = 0 $$
        where $\gamma_1 = \phi_1 + \psi_1$ is the perturbation to the gravity potential and $g_0 = d\phi_0/dr$.
    \end{alertblock}
    \pause
    \begin{block}{Boundary Condition}
        The condition that pressure is zero on the perturbed surface, $r = b+sh_1$, gives a linearised boundary condition on the reference surface $r=b$:
        $$ p_1 - \rho_0 g_0 h_1 = 0 \quad @ \ r=b $$
    \end{block}
\end{frame}

\begin{frame}{Slowly Rotating Planets (3/3): Sketch of the Solution}
    The linearised equation can be manipulated to relate all perturbations to the gravity potential perturbation, $\gamma_1$.
    \begin{itemize}
        \item Taking $\hat{\mathbf{r}} \times (\text{Eqn})$ shows that $\nabla(p_1 + \rho_0\gamma_1)$ is radial. Assuming perturbations average to zero on spheres, this gives:
        $$ p_1 = -\rho_0 \gamma_1 $$
        \pause
        \item Taking $\nabla \times (\text{Eqn})$ similarly shows that:
        $$ \rho_1 = \frac{1}{g_0} \frac{d\rho_0}{dr} \gamma_1 $$
        \pause
        \item The boundary condition then gives the shape perturbation:
        $$ h_1 = -\frac{\gamma_1}{g_0} \quad @ \ r=b $$
    \end{itemize}
    \begin{alertblock}{Main Result}
        All first-order perturbations ($p_1, \rho_1, h_1$) are directly proportional to the gravity potential perturbation, $\gamma_1 = \phi_1 + \psi_1$.
    \end{alertblock}
\end{frame}

\begin{frame}{Clairaut's Theory}
    \begin{itemize}
        \item We have now expressed all unknown perturbations in terms of the gravity potential perturbation, $\gamma_1$.
        \pause
        \item The final piece of physics is \textbf{Poisson's equation}, which relates the gravitational potential perturbation $\phi_1$ back to the density perturbation $\rho_1$ and shape perturbation $h_1$.
        $$ \nabla^2\phi_1 = 4\pi G \rho_1 $$
        \pause
        \item Substituting our expressions for $\rho_1$ and $h_1$ into Poisson's equation and its boundary conditions results in a solvable differential equation for $\phi_1$ in terms of the known centrifugal potential $\psi_1$.
    \end{itemize}
    \pause
    \begin{block}{Conclusion: Clairaut's Theory}
        This procedure, first developed by Clairaut, allows one to calculate the shape and internal structure of a slowly rotating planet in hydrostatic equilibrium. The resulting shape is an \textbf{oblate spheroid}, flattened at the poles. For the Earth, the equatorial radius is about 21 km larger than the polar radius.
    \end{block}
\end{frame}

% --- SUMMARY ---

\begin{frame}{Summary: What You Need to Know}
    \begin{itemize}
        \item A planet in \textbf{hydrostatic equilibrium} has its stress determined solely by pressure.
        \pause
        \item This requires that its surfaces of constant pressure, density, and gravity potential ($\gamma = \phi + \psi$) must all coincide.
        \pause
        \item A non-rotating hydrostatic planet must be spherically symmetric.
        \pause
        \item The shape of a \textbf{slowly rotating} planet can be found using perturbation theory. The key result is that all perturbed quantities (pressure, density, shape) are proportional to the gravity potential perturbation.
        \pause
        \item This leads to \textbf{Clairaut's theory}, which predicts the planet's shape to be an \textbf{oblate spheroid}.
    \end{itemize}
\end{frame}

\end{document}
